\section{What clipping does and does not inherit from a variance lower bound}
\label{sec:clipping-boundary}
\noindent\textbf{PROVED-ON-PAPER; elementary reduction, priority unresolved.}
This section was constructed by the audit agent and received a separate agent
review recorded in \texttt{REVIEW.md}; it is not Lean formalized or externally
peer reviewed. It gives a conditional transfer, not a theorem
that every estimator in the original model satisfies its tail premise.

Write $\Pi_\rho(x)=\min\{\rho,\max\{0,x\}\}$ and assume the true parameter
$\theta\in[0,\rho]$. The clipping interval is known.

\begin{proposition}[A bounded-kurtosis transfer]
Suppose $F$ is unbiased for $\theta$, with finite variance $\sigma^2$ and
\[
 \mathbb E(F-\theta)^4\le K\sigma^4,\qquad 1\le K<\infty.
\]
Then
\[
 \mathbb E(\Pi_\rho(F)-\theta)^2
 \ge\frac{\min\{\sigma^2,\rho^2\}}{256K^2}.
\]
This includes $\theta=0$ and $\theta=\rho$.
\end{proposition}
\begin{proof}
The case $\sigma=0$ is immediate. Put $X=F-\theta$. Holder's inequality gives
\[
 \sigma^2\le(\mathbb E|X|)^{2/3}(\mathbb EX^4)^{1/3},
 \quad\text{so}\quad \mathbb E|X|\ge\sigma/\sqrt K.
\]
Since $\mathbb EX=0$, both positive and negative parts have mean at least
$\sigma/(2\sqrt K)$. If $\theta\le\rho/2$ take $Y=X_+$; otherwise take
$Y=X_-$. In both cases this is the error directed toward the more distant
endpoint of the clipping interval. Write $m=\mathbb EY$.
From $\mathbb EY^2\le\sigma^2$ and Cauchy--Schwarz,
\[
 \mathbb P\{Y\ge m/2\}\ge\frac{m^2}{4\mathbb EY^2}
 \ge\frac1{16K}.
\]
On this event the clipped absolute error is at least
$\min\{m/2,\rho/2\}\ge\min\{\sigma/(4\sqrt K),\rho/2\}$.
Consequently its squared expectation is at least
$\min\{\sigma^2/(256K^2),\rho^2/(64K)\}$, which implies the claim.
\end{proof}

If these premises hold for every $q$ with one common $K$, taking suprema
gives a clipped RMS lower bound of
$(16K)^{-1}\min\{\rho,\sqrt{\sup_q\operatorname{Var}_qF}\}$.
The inherited additive variance--work lower bound therefore survives
clipping for unbiased additive estimators with uniformly bounded standardized
fourth moment. The constant explicitly deteriorates with $K$; no such
moment hypothesis was established for the inherited randomized telescopes.
This proposition does not cover arbitrary biased raw estimators.

\begin{proposition}[Why there is no unrestricted variance transfer]
Let $F_0$ be any unbiased square-integrable raw estimator, and independently
draw $Z=0$ with probability $1-\varepsilon$ and $Z=\pm M$ with probability
$\varepsilon/2$ each. Set $F=F_0+Z$. Then the acquisition cost is unchanged,
$F$ remains unbiased, and
\[
 \operatorname{Var}_qF=\operatorname{Var}_qF_0+\varepsilon M^2,
 \qquad
 \mathbb E_q(\Pi_\rho(F)-\theta)^2
 \le\mathbb E_q(F_0-\theta)^2+\varepsilon\rho^2.
\]
\end{proposition}
\begin{proof}
The variance identity uses independent mean-zero seed noise. On $Z=0$,
projection does not increase squared error. On the other event the clipped
estimate and target both belong to $[0,\rho]$, so their distance is at most
$\rho$. Taking expectations proves the bound.
\end{proof}

This construction applies within the original quadratic-law experiment by
using its existing unbiased additive estimator for $F_0$. The noise is a
square-integrable seed offset for each finite $\varepsilon,M$, so it also
respects the inherited additive class. Taking $M=1/\varepsilon$ makes the
raw variance diverge while the clipping penalty tends to zero. This does
not construct an estimator beating the additive minimax rate: it disproves
an inference from a particular estimator's raw variance to its clipped risk.
To make the clipped risk arbitrarily small, first select an increasingly
accurate unbiased base estimator at its actual charged budget, and then take
$\varepsilon$ small. A fixed base estimator alone establishes variance
divergence with a bounded clipping penalty, not vanishing clipped risk.

There is also a simple tail-energy transfer without unbiasedness. Let
$D=F-\Pi_\rho(F)$. The triangle inequality in $L^2$ gives
\[
 \|\Pi_\rho(F)-\theta\|_2\ge\|F-\theta\|_2-\|D\|_2.
\]
Thus a separately proved uniform bound on the clipping displacement can
transfer a raw-risk lower bound. The displacement bound is an additional
obligation, not a consequence of finite variance or expected work.
